\scanpage{81}
Localisation pour les variétés abéliennes.

1.\quad \inserted{Soit $A$ un anneau commutatif.} Soit $C$ une catégorie $A$-linéaire (pas nécessairement additive; les $\operatorname{Hom}$ sont munis de structures $A$-linéaires, les accouplements sont $A$- bilinéaires). Pour toute\deleted{\unsure{anneau $R$}} $A$-algèbre $A'$, soit $C_{A'}$ la catégorie ayant mêmes objets que $C$, et où les $\operatorname{Hom}$ sont
\[
\operatorname{Hom}_{C_{A'}}(X,Y)=\operatorname{Hom}_C(X,Y)\otimes_A A'.
\]
On trouve ainsi une\deleted{champ en}\inserted{catégorie fibrée \deleted{en}\inserted{\handwritten{\uline{$C$}}}} catégories $\unsure{\mathbb{R}|\mathbb{A}}$-linéaires sur le site annelé des schémas affines $\operatorname{Spec}(A')$ sur $S=\operatorname{Spec}(A)$. Elle n'est pratiquement jamais un champ sur ce gros site, même p\unsure{\deleted{i}o|\deleted{o}i}ur une topologie assez grossièr\unsure{\deleted{i}e|é|\deleted{e}i} comme la topologie de Zariski. Cependant, si on se restreint au petit site de Zariski \inserted{affine} de $S$ des ouverts affines de $S$, on a souvent un champ (pour la topologie de Zariski), \deleted{Ceci a des conséquences}\inserted{soit $\uline{C}$.} \handwritten{On dit alors que $C$ est \uline{localisable} \uline{sur} $A$.}
\margin{\handwritten{\uuline{$C$}}}

Supposons cette condition\deleted{s} satisfaite. On\inserted{\handwritten{n}} a des conséquences agréables, comme:

a)\quad So\unsure{\deleted{i}p|\deleted{p}i}t \deleted{$E$} $R$ une $A$-algèbre (pas néc. commutative), et $X$ un objet de $C$ muni d'une opération $\mathbb{R}$-linéaire de $R$. Alors pour tout $R$-module à d\unsure{\deleted{o}r|\deleted{r}o}oite $P$, qui est isomorphe localement sur $S$ à $R_{\unsure{d|a}}$, le produit tensoriel $P\otimes_R X$ dans $C$ existe, et \inserted{si $R\quad\operatorname{End}_C(X)$ est bijectif,} le foncteur
\[
P\qquad P\otimes_R X
\]
de la catégorie \inserted{$\uline{P}(R)$} des $R$-modules à droite localement isomorphes à $R$ sur $S$ dans la catégorie $C$ est pleinement fidèle, \handwritten{et son image essentielle \unsure{est} formée des $X'\in\operatorname{Ob}C$ qui sont ``localement isomorphes à $X$'' sur $A$.}

b)\quad Supposons que $C$ soit additive. Alors $\uline{C}$ est un champ en catégories additives $\mathbb{R}$-linéaires. Pour tout $X$, $R$ comme dessus, \deleted{et} et tout module à droite $P$ sur $R$ qui est localement sur $S$ libre de type fini, le produit tensoriel $\mathbb{P}\otimes_R X$ existe, et si $R\quad\operatorname{End}_C(X)$, le foncteur $\unsure{\deleted{P}|\mathbb{P}|\mathbb{R}}$
\[
P\qquad P\otimes_R X
\]
% archive sheet number: 80
